\documentclass[a4paper, 12pt]{report}

%-----------------------------------------------------
%                      PACKAGES
% ----------------------------------------------------
\usepackage[T1]{fontenc}
\usepackage[latin9]{inputenc}
\usepackage{fancyhdr}
\usepackage{geometry}
\usepackage[colorlinks=true, linkcolor=blue,citecolor=blue, urlcolor=blue]{hyperref}
\usepackage{indentfirst}
\usepackage{graphicx}
\usepackage{float}
\usepackage{setspace}
\usepackage{amsmath}
\usepackage{amssymb}
\usepackage{multirow}
\usepackage[table,xcdraw]{xcolor}
\usepackage{colortbl}
\usepackage{tabularx,booktabs}
\usepackage{enumerate}
\usepackage{titlesec}
\usepackage{caption}
\usepackage{subcaption}
\usepackage{nicefrac, xfrac}
\usepackage{subfiles}
\usepackage{mathtools}
\usepackage{bm}
\usepackage{tikz}
\usepackage[mode=buildnew]{standalone}
\usepackage{systeme}
\usepackage{stmaryrd}
\usepackage{import}
\usepackage[english]{babel} 
\usepackage{times}
\usepackage{pgfplots}
\usepackage{makecell}
\usepackage{threeparttable}
\usepackage{ragged2e}
\usepackage{tocloft}
\usepackage{booktabs}
\usepackage[acronym]{glossaries}
\usepackage{pdfpages}
\usepackage{lipsum}
\usepackage{hyphenat} 
\usepackage{xspace}
\usepackage{abntex2cite}

%-----------------------------------------------------
%               DOCUMENT SETTINGS
% ----------------------------------------------------
\geometry{a4paper,top=30mm,bottom=20mm,left=30mm,right=20mm}

\pgfplotsset{compat=1.17} 

% Defining spaces between lines
\setstretch{1.5} 

\graphicspath{{Figures/}}

\tikzset{
	font={\fontsize{11pt}{12}\selectfont}}

\setcounter{secnumdepth}{4}
\setcounter{tocdepth}{4}

% defining the height of table cells
\renewcommand{\arraystretch}{1.2}

\def\textcite#1{\citeauthoronline{#1} \cite{#1}}
%-----------------------------------------------------
%              PRE TEXTUAL ELEMENTS
% ----------------------------------------------------
\newenvironment{epigrafe}{\newpage\mbox{}\vfill\hfill\begin{minipage}[t]{0.5\textwidth}}
{\end{minipage}\newpage}

%-----------------------------------------------------
%               GLOSSARY & ACRONYMS
% ----------------------------------------------------
% theorically, you should have a file for each glossary and acronym but for the sake of simplicity, I will put everything in the same file here so every chapter can access it
\makeglossary

\newacronym{cfd}{CFD}{Computational Fluid Dynamics}

\newglossaryentry{a1}{
    name = {A1},
    description = {This is the description of A1}
}

%-----------------------------------------------------
%                     DOCUMENT 
% ----------------------------------------------------
\begin{document}
% ------- WORK INFORMATION -------
\author{Carlos Henrique Chama Puga}
\title{Numerical Simulation of Head Loss in Pipes for Oil Extraction}

\def\titulo{Simula\c{c}\~{a}o Num\'{e}rica de Perda de Carga em Tubula\c{c}\~{o}es para Extra\c{c}\~{a}o de Petr\'{o}leo}

\def\Uni{{Universidade Estadual de Campinas}}
\def\Fac{{Faculdade de Engenharia Civil, Arquitetura e Urbanismo}}

\def\advisor{Prof. Dr. Philippe Remy Bernard Devloo}
\def\coadvisor{Prof. Dr. Nathan Shauer}

\makeatletter

% ------- PRE TEXTUAL PAGES -------
% Could not find another way to remove the page number from the first pages
\fancypagestyle{plain}{
    \fancyhf{}% Limpa todos os campos
    \fancyfoot[C]{}%
    \renewcommand{\headrulewidth}{0pt}%
}

\pagestyle{plain}
\def\logos{
    \noindent
    \raisebox{-.5\height}{\includegraphics[width=2.2cm]{Figures/logo-unicamp.pdf}}

    \vspace*{2cm}
    
    \noindent
    \begin{center} \large
        \MakeUppercase{\Uni}\\
        \Fac\\
    \end{center}
}

\def\openningpage{
  \logos
  \vskip 35mm
  \begin{center}
    \Large
    {\bf \@author}
    \vskip 25mm
      {\bf \@title}
      \vskip 25mm
    {\bf \titulo}
    \vfill
    \large
    Campinas\\\workyear
  \end{center}
}

\openningpage % Cover page

\newpage
\def\titlepage{
    \large\bf
    \vfill
    \begin{center}
        \@author
        \vfill
            \@title
            \vskip 8mm
        \titulo
    \end{center}
    \vfill
    \normalfont
    \normalsize
    \hfill\begin{minipage}[t]{0.54\textwidth}
        Disserta\c{c}\~{a}o de Mestrado apresentada \`{a} Faculdade de Engenharia Civil, Arquitetura e Urbanismo da Universidade Estadual de Campinas como parte dos requisitos para a obten\c{c}\~{a}o do t\'{i}tulo de Mestre, na área de Engenharia Civil.
    \end{minipage}
        \vskip 8mm
        \hfill\begin{minipage}[t]{0.54\textwidth}
            Master's Degree Thesis presented to the School of Civil Engineering, Architecture and Urban Planning ofthe Universidade Estadual de Campinas in partial fulfillment of the requirements for the degree of Master of Science in Civil Engineering.
        \end{minipage}
        \vskip 8mm
        \vfill
        {\bf Advisor/}\hspace{-0.9ex}
        {\bf Orientador: \advisor}

        % comment this line if you don't have a co-advisor
        {\bf Co-advisor/}\hspace{-0.7ex}
        {\bf Co-orientador: \coadvisor}
    \vfill
        \begin{minipage}[t]{0.5\textwidth}
            Este exemplar corresponde \`{a} vers\~{a}o final da Disserta\c{c}\~{a}o de Mestrado apresentada por \@author, orientada pelo \advisor \xspace, e co-orientada pelo \coadvisor.
        \end{minipage}
    \vfill
    \begin{center}
        \large
        Campinas\\ 2024
    \end{center}
}

\titlepage % Description page

% when you already have the catalografi card, you need to comment the next line
\newpage
\begin{center}
    \topskip0pt
    \vspace*{\fill}
    \textit{This page is intentionally left blank.}  
    \vspace*{\fill}
\end{center}
 % Blank page
% and add the Catalografic card
% \\includepdf[pages=-]{FirstPages/CatalograficCard.pdf} 

% if you don't want to include the epigraph, you can comment the next line
\newpage
\begin{epigrafe}
    \thispagestyle{empty}
    \centering
    \textit{``Quote here.''}
    \begin{flushright}
       - Author
    \end{flushright}
\end{epigrafe} % Epigraph

% including the aknowledgements is a good practice
\chapter*{Acknowledgements}
Here you can write your acknowledgements. This is a very important part of your work, so take your time to write it. You can thank your family, friends, colleagues, professors, and anyone else who helped you during your research. You can also thank the funding agencies that supported your work. % Acknowledgements

% every work should have an abstract (in Brazil, you should have a resumo too)
\chapter*{Resumo}
    Este \'{e} o resumo em portugu\^{e}s do trabalho.


\chapter*{Abstract}
    This is the abstract in English of the work. % Abstract

% List of figures
\newpage
\listoffigures 

% List of tables
\newpage
\listoftables 

\printglossary[nonumberlist] % Glossary
\printglossary[type=\acronymtype, nonumberlist] % Acronyms

% Table of contents
\newpage
\tableofcontents 

% from here on, the page number is shown
% Chapter first page settings
\newpage
\fancypagestyle{plain}{
    \fancyhf{}% Limpa todos os campos
    \fancyhead[R]{\thepage}%
    \renewcommand{\headrulewidth}{0pt}%
}

\fancypagestyle{headings}{%
    \fancyhf{}% Limpa todos os campos
    \fancyhead[L]{\textsc{\nouppercase{\leftmark}}}
    \fancyhead[R]{\thepage}% Numero da página à direita
    \renewcommand{\headrulewidth}{1pt}%
}

% ------- CHAPTERS -------
\pagestyle{headings}
\chapter{sjnsa}
This is something
\begin{equation}
    a = 1,
\end{equation} \label{eq:1}
\autoref{eq:1} represents something

\newpage
\section{aans}
\cite{Lima_Thesis}

% ------- BIBLIOGRAPHY -------
\addcontentsline{toc}{chapter}{Bibliography}
\bibliographystyle{abntex2-num}
\bibliography{References}

% ------- APPENDIX -------
\appendix
\chapter{Latex Questions}
If you have any questions about the \LaTeX, check \href{https://www.latex-project.org/}{here} for more information. 
\end{document}